\section{Background: amplification, audits and platform governance}\label{sec:background}

This section reviews what published audits of recommender systems have done and found, the methodological debates they have opened, the regulatory setting in which an audit of a large service now takes place, and the problem of measuring antisemitic content with a classifier. Three conventions apply. First, every empirical finding below is reported as its authors state it, with the scope they give it: the service, the period, the design and the content category. We do not extend any finding beyond that scope. Second, almost none of the studies measures antisemitic content as such. Most study political ideology, ``problematic'' or extremist channel categories, or incivility and slurs, so their findings bear on our question by analogy only. Third, none of these findings is an input to the simulation study of \cref{sec:simulation}, and we have run no audit of any real service.

\subsection{Algorithm audits and what they have found}\label{sec:bg-audits}

\paragraph{Audit designs}
\Citet{sandvig2014} set out the research designs available to outsiders who wish to learn how a platform's algorithm treats different users or content when its code and data are not available. They range from inspecting code, through surveying users, scraping and crowdsourcing data from real users, to \emph{sock-puppet} audits, in which researchers create accounts that behave in controlled ways and record what the system shows them. The sock-puppet design is the one closest to an experiment. The researcher fixes the account's behaviour and observes the system's response. It is the template for the counterfactual designs of \cref{sec:identification}.

\paragraph{Non-personalised crawls}
\Citet{ribeiro2020} audited radicalisation pathways on YouTube. They analysed 330,925 videos from 349 channels, which they classified into four groups: Media, the Alt-lite, the Intellectual Dark Web (I.D.W.) and the Alt-right. Using user comments as a trace of consumption, they report that the three non-media communities increasingly share a commenting user base and that users migrate from milder to more extreme content. They also collected more than two million video and channel recommendations between May and July 2019, and report that Alt-lite content is easily reachable from I.D.W.\ channels, while Alt-right videos are reachable only through channel recommendations. The authors state that the recommendation data do not include personalisation, because no account was logged in, and that personalisation could reveal a different picture. \Citet{ledwich2020} also collected recommendations without logging in, across a set of political channels. They conclude that the recommendation algorithm did not promote radicalising or inflammatory content and instead favoured mainstream media channels. Because their recommendations were not personalised, critics argued that the design could not speak to the ``rabbit hole'' of an individual viewer with a history. The two studies thus reached different conclusions from data of a similar, non-personalised kind, and both concern a single service in 2019.

\paragraph{Sock puppets with trained histories}
\Citet{whittaker2021} ran an experimental audit on three platforms, YouTube, Reddit and Gab, designed to account for personalisation. They report that YouTube amplified extreme and fringe content, while Reddit and Gab did not. \Citet{haroon2023} trained 100,000 sock puppets on YouTube to have ideological watch histories and then followed the ``up next'' trail of recommendations. They report that the algorithm recommends ideologically congenial content to partisan puppets, with some moderate and cross-cutting exposure, and that congenial recommendations increase deeper in the trail for right-leaning puppets. They did not find meaningful increases in the ideological extremity of recommendations deeper in the trail. They did find that a growing proportion of recommendations came from channels they categorised as problematic (for example ``IDW'', ``Alt-right'', ``Conspiracy'' and ``QAnon''), that the increase was most pronounced for very-right puppets, and that the overall proportion of such recommendations was low. A non-peer-reviewed report by ADL and the Tech Transparency Project set up six test personas of different ages and genders on four platforms. It reports that three of the four recommended explicitly antisemitic and extremist content to the personas, and that personas which followed or liked the suggested accounts were shown more such content \citep{adl2023badtoworse}. That report is one of the few audits that targets antisemitic content directly. With six personas and no reference policy, however, it shows that such recommendations occurred, not how much more often they occurred than under an alternative policy.

\paragraph{Counterfactual bots}
\Citet{hosseinmardi2024} addressed the missing counterfactual directly. Their bots were first primed with the histories of real users. In the observation phase, some bots continued to replicate the real users' consumption, while ``counterfactual'' bots followed rule-based trajectories through the recommendations. They report that relying exclusively on the recommender results, on average, in less partisan consumption, with the effect most pronounced for heavy partisan consumers. In a second experiment, when bots switched from partisan to moderate content, the sidebar recommender ``forgot'' their partisan preference within roughly 30 videos regardless of prior history, while homepage recommendations shifted more gradually. The authors conclude that, at least since the algorithm changes the platform implemented in 2019, individual consumption mostly reflects individual preferences, with recommendations playing, if anything, a moderating role. The qualifier ``at least since 2019'' is the authors' own, and it limits the finding in time.

\paragraph{Real-user data}
\Citet{hosseinmardi2021} analysed the browsing histories of a panel of more than 300,000 Americans, on and off YouTube, from January 2016 to December 2019. They report little evidence that the recommendation algorithm drives attention to radical political content, either left or right, and find that consumption of such content tracks users' broader preferences across the web. \Citet{chen2023} paired survey and browser-activity data from 1,181 participants recruited from a representative sample in 2020. They report that exposure to videos from alternative and extremist channels was concentrated among a small group with high prior levels of gender and racial resentment, who typically subscribed to these channels and also arrived through external links, and that non-subscribers rarely saw or followed recommendations to such videos. They suggest that the algorithm was not sending people down ``rabbit holes'' during their 2020 observation window, possibly because of the 2019 changes, while the platform still played a key role in exposing a dedicated audience to this content. Outside video, \citet{robertson2023} compared the news sources that participants were shown in Google Search results with the sources they engaged with. They report that engagement included more identity-congruent and unreliable sources than exposure did, and conclude that exposure to and engagement with partisan or unreliable news on that service were driven mainly by users' own choices rather than by algorithmic curation.

\paragraph{Randomised holdouts and feed experiments}
Two studies randomised the feed policy itself. \Citet{huszar2022}, most of whose authors were affiliated with the platform, used a long-running randomised experiment on Twitter in which a control group of nearly two million daily active accounts received a reverse-chronological home timeline without algorithmic personalisation. They compared the reach of tweets from elected legislators in seven countries, and of United States news sources, between the personalised and control groups. They report that in six of the seven countries the mainstream political right enjoyed higher algorithmic amplification than the mainstream political left, that algorithmic amplification favoured right-leaning news sources, and that they found no evidence for the hypothesis that the algorithm amplifies far-left and far-right groups more than moderate ones. \Citet{guess2023} randomly assigned consenting users of Facebook and Instagram to a reverse-chronological feed for three months during the 2020 United States election. They report that the chronological feed substantially reduced time spent and activity on the platforms; increased the amount of political and untrustworthy content users saw on both platforms; decreased, on Facebook, the amount of content classified as uncivil or as containing slur words; and did not significantly alter issue polarisation, affective polarisation, political knowledge or other key attitudes. It is a close published analogue of our reference-policy contrast, with an outcome adjacent to hateful content. Its outcome was measured by classifiers, and its later erratum reports that the incivility and slur classifiers were not applied to one post format on one platform, so the reported proportions for that format were likely underestimated. Classifier coverage, in other words, was part of the result.

\begin{table}[tbp]
\centering
\caption{Published studies of recommendation and exposure discussed in \cref{sec:bg-audits}, as reported by their authors. Entries summarise design and scope only; see the text for each study's own statement of its findings. None measures antisemitic content with a validated classifier, and only \citet{adl2023badtoworse} targets antisemitic content.}\label{tab:bg-studies}
\footnotesize
\setlength{\tabcolsep}{4pt}
\renewcommand{\arraystretch}{1.15}
\begin{tabularx}{\textwidth}{@{}>{\raggedright\arraybackslash}p{2.6cm}>{\raggedright\arraybackslash}p{2.0cm}>{\raggedright\arraybackslash}X>{\raggedright\arraybackslash}p{3.6cm}@{}}
\toprule
Study & Service(s) & Design and comparison & Content measure \\
\midrule
\citet{ribeiro2020} & YouTube & Logged-out recommendation crawl (May--Jul 2019); comment trajectories & Channel categories (Media, Alt-lite, I.D.W., Alt-right) \\
\citet{ledwich2020} & YouTube & Logged-out recommendation crawl & Channel categories \\
\citet{whittaker2021} & YouTube, Reddit, Gab & Experimental accounts with personalisation & Extreme and fringe content \\
\citet{haroon2023} & YouTube & 100,000 trained sock puppets; up-next trails & Ideology; ``problematic'' channel list \\
\citet{adl2023badtoworse} & Four platforms & Six test personas; no reference policy (NGO report) & Antisemitic and extremist content \\
\citet{hosseinmardi2024} & YouTube & Counterfactual bots primed with real histories & Partisanship of content \\
\citet{hosseinmardi2021} & YouTube & Real-user browsing panel, 2016--2019 & Channel categories (incl.\ far right) \\
\citet{chen2023} & YouTube & Survey plus browser data, $n=1{,}181$, 2020 & Alternative and extremist channels \\
\citet{robertson2023} & Google Search & Exposure versus engagement of participants & Partisan and unreliable news sources \\
\citet{huszar2022} & Twitter & Randomised chronological holdout & Reach of legislators and news sources \\
\citet{guess2023} & Facebook, Instagram & Randomised assignment to chronological feed, 3 months, 2020 & Political, untrustworthy, uncivil and slur content (classifiers) \\
\bottomrule
\end{tabularx}
\end{table}

\paragraph{What the literature does and does not establish}
\Cref{tab:bg-studies} summarises the designs. The findings vary by service, by period and by design. Several studies of one video service report a weaker algorithmic push, or none, towards partisan or extreme content after 2019, and attribute this, tentatively, to changes the service made that year \citep{hosseinmardi2024,chen2023}. Other designs on the same service, and studies of other services, report congenial or problematic recommendations for some profiles \citep{whittaker2021,haroon2023}. The randomised feed experiments report amplification that differs by political group \citep{huszar2022}, and a chronological alternative that changes the mix of content users see in both directions \citep{guess2023}. These results are not in conflict with one another so much as answers to different questions. They differ in the estimand (what is recommended, what is consumed, or what users come to believe), in the comparison (no comparison, a logged-out crawl, counterfactual bots, or a randomised chronological feed), in the content category, and in the period. A single audit cannot settle whether ``recommender systems amplify antisemitism''. It can estimate a well-defined contrast for a stated service, surface, period and demand profile. \Cref{sec:estimands} defines such contrasts.

\paragraph{Feedback loops}
Audits observe a system that learns from the behaviour it induces. \Citet{chaney2018} show in simulation that training recommenders on data generated under earlier recommendations increases the homogeneity of user behaviour and decreases utility. \Citet{jiang2019} study theoretically how the feedback between a recommender and the evolving interests of its users can degenerate, so that what the system shows and what the user prefers reinforce one another. For an audit, these loops have two consequences. Exposure under the algorithmic policy depends on the persona's own history, which is why \cref{sec:markov} models a session as a Markov chain. And puppets that feed the same learning system can change what it recommends to one another, which is the interference problem of \cref{prop:interference}.

\subsection{Methodological debates}\label{sec:bg-debates}

\paragraph{Sock puppets or real users}
Real-user data show what people actually saw and chose, but they cannot by themselves separate the algorithm's contribution from the user's intentions, because users who want extreme content both seek it and receive recommendations for it. \Citet{hosseinmardi2024} motivate their counterfactual bots by exactly this missing counterfactual. Sock puppets supply the counterfactual, since the researcher fixes behaviour and varies the policy. The price is realism. A puppet's demand is a script, so its exposure is a property of the policy evaluated at that script, not a description of any real population. Our estimands embrace this. The amplification ratio of \cref{def:amplification} is defined \emph{for a fixed demand profile} $d$, so it is a statement about supply conditional on demand. Population-level statements require weights over demand profiles, which an audit does not provide.

\paragraph{Logged out or logged in}
Logged-out crawls \citep{ribeiro2020,ledwich2020} measure a non-personalised recommender that depends only on context, such as the current item. They impose the smallest footprint and need no account. However, they cannot measure personalisation, which is the mechanism most often suspected of amplification. Logged-in puppets with trained histories \citep{whittaker2021,haroon2023} can measure personalisation, but they require accounts and therefore engage the services' terms (\cref{sec:ethics}). In our framework, a logged-out view is not a flawed version of the personalised audit. It is one possible \emph{reference policy}: the contrast between a trained puppet's feed and a logged-out session that follows the same choice rule isolates the contribution of profile-based personalisation (\cref{sec:protocol}).

\paragraph{Demand or supply}
\Citet{munger2022} propose a supply-and-demand framework for political video. They argue that the consumption of white-nationalist video was not caused by the supply of such media radicalising an otherwise moderate audience. The real-user studies above \citep{hosseinmardi2021,chen2023,robertson2023} likewise report that consumption of extreme or unreliable content tracks prior preferences and choices. These are claims about which side of the market accounts for most observed consumption. They are compatible with an algorithm that supplies more antisemitic content to some demand profiles than a reference policy would, and also with one that supplies less. The decomposition of \cref{thm:decomp} separates observed exposure into a demand-driven part, an algorithm-driven part and an explicit interaction term, so that the supply-or-demand question becomes an estimable contrast rather than a dichotomy.

\subsection{Platform governance and research access}\label{sec:bg-governance}

\paragraph{Systemic-risk assessment under the Digital Services Act}
The EU Digital Services Act (DSA), Regulation (EU) 2022/2065, imposes additional obligations on providers of very large online platforms and very large online search engines \citep{eu2022dsa}. Article~34 requires them to identify, analyse and assess systemic risks stemming from their services, including the dissemination of illegal content and negative effects on the exercise of fundamental rights. Article~34(2) lists the design of their recommender systems and any other relevant algorithmic system among the factors whose influence on those risks the assessment must take into account. Article~35 requires reasonable, proportionate and effective measures to mitigate the risks identified, and the recitals name testing and, where necessary, adapting algorithmic systems, including recommender systems, among such measures (Recital~88). Article~38 requires these providers to offer, for each of their recommender systems, at least one option that is not based on profiling as defined in the General Data Protection Regulation. Where such an option exists, it is a candidate reference policy that the service itself supplies (\cref{sec:protocol}). Whether a particular non-profiling option corresponds to a reverse-chronological feed, a popularity ranking or something else is a matter of fact about each service; the Regulation does not prescribe it.

\paragraph{Data access for researchers}
Article~40 of the DSA provides two routes to data. Under Article~40(4), upon a reasoned request from the Digital Services Coordinator of establishment, providers must give vetted researchers access to data for research that contributes to the detection, identification and understanding of systemic risks in the Union and to the assessment of risk-mitigation measures. Under Article~40(12), providers must give researchers who meet the stated conditions, including independence from commercial interests and disclosure of funding, and who use the data for research on systemic risks, access without undue delay to data that is publicly accessible in their online interface, including real-time data where technically possible \citep{eu2022dsa}. The technical conditions and procedures for the vetted route are set by Commission Delegated Regulation (EU) 2025/2050 of 1 July 2025, published in the Official Journal on 9 October 2025 \citep{eu2025delegated}. The Commission announced that the new rules applied from 29 October 2025, with applications made through a DSA data access portal \citep{ec2025access}. As of 19 May 2026, according to the Commission, Digital Services Coordinators had received 49 applications for assessment, mainly requesting data from social media platforms, and a roundtable with providers had been held to prepare the first data requests \citep{ec2026roundtable}. We have found no published account of how many applications had been granted by the time of writing. The route is therefore operational in law, but its practical reach for a given project, service and period is still being established. Separately, in its first non-compliance decision under the DSA, adopted in December 2025, the Commission found among other things that one provider's terms prohibited eligible researchers from independently accessing its public data, including by scraping \citep{ec2025fine}.

\paragraph{Transparency reports}
The DSA also requires periodic transparency reports. Providers of very large online platforms must publish them at least every six months (Article~42), using templates and harmonised reporting periods (January to June and July to December) laid down in Commission Implementing Regulation (EU) 2024/2835 \citep{eu2024transparency}. The results of risk assessments, the mitigation measures adopted and the reports of independent audits are also to be made public, at the latest three months after the receipt of each audit report (Article~42(4)). These documents report moderation activity and risk assessments at the level of the service. They do not report exposure by demand profile under alternative policies, which is the quantity an amplification audit estimates. Independent audits of the kind proposed here complement them rather than duplicate them.

\subsection{Classifier-based measurement of antisemitic content}\label{sec:bg-classifier}

Every quantity in this paper is defined in terms of the true label $L_c$ of each item, but an audit observes only a classifier's label $\widehat L_c$. The lab's detection proposal frames the evaluation of such classifiers as a contextual problem: classifiers must distinguish antisemitic expression from quotation, criticism and other contextual uses, and their errors should be reported through precision and recall with uncertainty rather than a single score \citep{kornsteinlab2026detection}. The companion measurement proposal adds that an estimate concerns the sampled corpus, not platform-wide prevalence \citep{kornsteinlab2026measurement}. Three difficulties matter for an audit.

\paragraph{Coded and evolving language}
Antisemitic content often uses coded terms whose meaning is not visible to a lexicon. \Citet{zannettou2020} studied large volumes of posts from two fringe web communities and used word embeddings to discover and categorise antisemitic terminology, including new coded terms, as it emerged. A classifier trained before an audit may therefore miss terms that emerge during it. Its sensitivity is a function of time and of the content mix, not a constant.

\paragraph{Context, quotation and counter-speech}
The same words can endorse antisemitism, report it, quote it in order to condemn it, or reclaim a term. HateCheck \citep{rottger2021} built functional tests for hate-speech detectors, including non-hateful cases in which hateful language is quoted or referred to in order to denounce it, and its authors report critical weaknesses in the models they tested. For an audit this matters directly. A policy that recommends news coverage of antisemitic incidents, or counter-speech, raises the share of items that quote antisemitic content without endorsing it. If the classifier confuses quotation with endorsement, the naive amplification ratio is biased even when the classifier's error rates are identical for each type of content, because the arms contain different mixtures of types. This is the differential-misclassification problem of \cref{thm:misclass}.

\paragraph{Definitions and annotation}
Ground truth is itself a coding decision. Annotated datasets of antisemitic content use explicit definitions. For example, \citet{jikeli2023} describe an annotation guide based on the IHRA Working Definition \citep{ihra2016}, under which each item was labelled by two annotators and items on which they disagreed were excluded from the final dataset. Excluding disagreements yields cleaner training data, but it removes exactly the borderline items that an audit's validation sample must cover. Our protocol therefore adjudicates disagreements rather than dropping them, and it reports scope variants under alternative definitions \citep{ihra2016,jda2021} side by side, as the companion paper does \citep{kornsteinlab2026darkfigure}. The validation sample, codebook and correction are specified in \cref{sec:protocol} and \cref{sec:measurement}.
